\chapter{Introduction}%
\label{ch:intro}

Top-quark pair (\ttbar) production is one of the most studied processes
at the Large Hadron Collider (LHC) and the measurement of differential cross
sections is of fundamental importance, since \ttbar\ production provides access
to fundamental parameters of the Standard Model (SM) of particle physics,
like the top-quark mass, but it is also an important background for many
other measurements and searches for physics Beyond the Standard Model (BSM).

The top-quark pair decays predominantly into two \Wboson bosons and two \bquarks (\WbWb).
Differential cross sections for \ttbar production have been
measured~\cite{TOPQ-2016-01,TOPQ-2018-15,CMS-TOP-16-007,CMS-TOP-16-014,CMS-TOP-17-002}
and calculated~\cite{Czakon:2011xx,Czakon:2015owf,Czakon:2016dgf} in a wide
kinematic range.
All these measurements consider the \WbWb\ final state.

Further processes contribute to the \WbWb final state, most prominently
the associated single-resonant top-quark production process \( tWb \)
(shown on a Feynman diagram in \cref{fig:intro:feynman}).
These contributions are considered as background
processes in top-quark pair production analyses, but these NLO real corrections
to the LO \tW production process interfere with \ttbar and with each other
and therefore cannot be unambiguously treated theoretically.
The interference between \ttbar and \tW therefore reduces the achievable precision for top-quark phenomenology
using top-quark pair cross-sections. In contrast, the direct measurement
of \WbWb production cross-sections provides an alternative access to top-quark
phenomenology, where those aspects are treated consistently in theory,
albeit at the cost of considerably more complicated calculations.
The interference between top-quark pair production and single-top production
at higher-order perturbative QCD also introduces a considerable uncertainty
in searches for new physics phenomena at large masses, and this
measurement aims to reduce this uncertainty for future searches.

\begin{figure}[htbp]
  \centering
  \includegraphics[width=\textwidth]{Theory/diagrams}
  \caption{Feynman diagrams for LO \ttbar (left) and NLO \(tWb\) (right) production processes.}%
  \label{fig:intro:feynman}
\end{figure}

This interference problem has been widely studied from the theoretical point
of view~\cite{White:2009yt,Frixione:2008yi,Hollik:2012rc,Demartin:2016axk,ATL-PHYS-PUB-2016-020}
and provides a significant source of uncertainty for many
measurements~\cite{TOPQ-2018-26,TOPQ-2020-06,TOPQ-2018-15}
and BSM searches~\cite{SUSY-2015-02,SUSY-2016-17,SUSY-2014-07,SUSY-2015-01,SUSY-2016-15,SUSY-2015-08}.
The consistent calculation of the \WbWb final state was performed up to next-to-leading order in QCD~\cite{Denner:2010jp},
including full top off-shell effects at NLO,
and supplemented with parton showers~\cite{Jezo:2016ujg,Jezo:2023rht}.
A calculation of the process \(pp\to\ell\nu b b j j\) was performed up to NLO in Ref.~\cite{Denner:2017kzu}.
The only existing measurement available to assess
these predictions in a region sensitive to interference is the one provided
by ATLAS in 2018~\cite{TOPQ-2017-05} in the dilepton channel.


In this note, we expand on that result to cover also semileptonic \WbWb decays
yielding a final state with a charged lepton, significant missing transverse momentum, two \bjets and light jets.
When compared to the dilepton channel, the semileptonic decay channel provides the opportunity to measure the \Wboson 
boson kinematics directly, since the hadronically decaying \Wboson boson can be fully reconstructed, and in addition, 
the larger cross section enables us to probe a larger phase space up to higher scales, which is important to constrain 
the \ttbar and \tW interference effects at high scales, which is relevant for BSM searches.

The analysis therefore focusses on three kinematic regions: a large phase space to constrain the interference effects and 
improve the modelling over a large kinematic region; an analysis region targeting the interference in phase-spaces typically 
addressed by BSM physics searches; and a dedicated measurement region with a focus on the \Wboson boson and \bjet kinematic observables. 
Differential cross-section measurements are performed for selected observables relevant for \ttbar physics.
These differential cross-sections can be used to assess and tune the performance
of new Monte Carlo generators, determine SM parameters, or constrain uncertainties in BSM searches.
The measurements are required to ensure that basic kinematics in the event
are well modelled when the interference is correctly modelled.
The full dataset collected by the ATLAS detector between 2015 and 2018
at \( \sqrt{s} = \qty{13}{\TeV} \) center-of-mass energy and
\( \mathcal{L} = \qty{140}{\ifb} \) is used for the measurement.

